%======================================================================%
\section{Introduction}
\label{sec:intro}
%======================================================================%

A unified description of gravitation and the Standard Model (SM) requires more
than phenomenological fit: the framework must be \emph{mathematically closed}
(definitions, axioms, and proofs form a coherent deductive system),
\emph{internally consistent} (anomalies cancel; representation dimensions and
$\beta$-function coefficients agree), and \emph{verifiable} (predictions and
falsification criteria are stated explicitly, not left to narrative).

\emph{Stone--Jordan Exceptional Theory} (SJET) is proposed as such a framework.
It is organized in three tiers:
\begin{enumerate}[label=\textbf{Tier \Roman*}.,leftmargin=*]
  \item \textbf{Discrete.} Configuration space is a weighted graph
    $\mathcal{G}$ with Boolean algebra $\mathcal{B}(\mathcal{G})$ and Stone
    spectrum $S=\mathrm{Spec}(\mathcal{B}(\mathcal{G}))$.
  \item \textbf{Arithmetic.} A profinite refinement tower carries the Hurwitz
    ladder $0\to\R\to\C\to\Quat\to\Oct$ and an exceptional Lie cascade terminating
    at $\jalg$.
  \item \textbf{Physical.} The order parameter $\orderparam\in\jalg$ breaks
    compact $\E{6}$ to $H=\Spin(10)\times U(1)$; soldering (Postulate~P1) yields
    four-dimensional Lorentzian spacetime; matter postulates complete SM descent.
\end{enumerate}

\subsection{Honest thesis}
\label{sec:intro:thesis}

SJET is an \textbf{$\E{6}$/Albert-algebra--motivated conditional embedding}.
Three axioms fix the arena and the capacity dynamics; postulates P1--P4 supply
the remaining structure needed for contact with four-dimensional quantum field
theory. Claims are never left implicit:
\begin{itemize}
  \item \textbf{Theorem} --- proved from stated hypotheses in this manuscript.
  \item \textbf{Postulate} --- added structure, named once near the axioms.
  \item \textbf{Conjecture} --- motivated but not controlled (e.g.\ gravitational
    UV fixed point, heterotic correspondence).
\end{itemize}
The phrase ``everything follows from three axioms'' is \emph{not} used. What
\emph{is} proved from the axioms alone includes the potential identification,
EIII vacuum selection, composite gauge kinematics, and anomaly cancellation
\emph{given} the $\mathbf{27}$ branching. SM fermion masses, three generations,
and Lorentzian signature require P1--P4.

\subsection{Relation to prior work}
\label{sec:intro:related}

SJET synthesizes several established lines: Hurwitz's classification of normed
division algebras~\cite{Hurwitz1898}; Freudenthal's exceptional geometry and the
Albert algebra~\cite{Freudenthal1954,Jacobson1968}; Stone duality~\cite{Stone1936};
capacity/entropy duality in optimal transport; $\E{6}$ coset models and
$\Spin(10)$ grand unification~\cite{GeorgiGlashow1974}; Green--Schwarz anomaly
cancellation~\cite{GreenSchwarz1984}; Sakharov induced gravity~\cite{Sakharov1967};
and asymptotic-safety ideas for gravity~\cite{Reuter1998}. The graph-to-algebra
lift and capacity-induced potential are the novel structural moves; their
convergence to the sharp-map potential \eqref{eq:dynamics:V} is
Theorem~\ref{thm:dynamics:potential}.

\subsection{Conventions}
\label{sec:intro:conventions}

\begin{itemize}[leftmargin=*]
  \item \textbf{Signature.} Spacetime metric $\eta_{ab}=\mathrm{diag}(-1,+1,+1,+1)$.
  \item \textbf{Real form.} Dynamics use compact $\E{6}$ (equivalently $E_{6(-14)}$,
    maximal compact $H=\Spin(10)\times U(1)$). Split $E_{6(-26)}$ classifies real
    orbits in the $\mathbf{27}$ only.
  \item \textbf{Order parameter.} $\orderparam$ denotes the $\jalg$-valued field;
    $\capacity$ denotes the graph capacity functional
    (Definition~\ref{def:ladder:capacity}). No symbol is overloaded.
  \item \textbf{Scales.} Decay constant $f$ has $[f]=\mathrm{mass}$; breaking scale
    $\mu$, UV cutoff $\Lambda$, GUT scale $\MGUT$; dimensionless couplings
    $\hat g^2$, $\xi$, $\hat y^2$, $\kappa$, $\Lambda_{\mathrm{cc}}$.
  \item \textbf{Casimirs.} $C_A(\Spin(10))=8$, $T(\mathbf{16})=2$ (index),
    $C_2(\mathbf{16})=45/8$; these are used once in all $\beta$-function sums.
  \item \textbf{Branching.} Adjoint:
    $\mathbf{78}=\mathbf{45}_0\oplus\mathbf{1}_0\oplus\mathbf{16}_{-3}
    \oplus\overline{\mathbf{16}}_{+3}$; matter:
    $\mathbf{27}=\mathbf{1}_{+4}\oplus\mathbf{10}_{-2}\oplus\mathbf{16}_{+1}$;
    coset tangent $\mathfrak{m}=\mathbf{16}_{-3}\oplus\overline{\mathbf{16}}_{+3}$.
\end{itemize}

\subsection{Organization}

Section~\ref{sec:foundations} states the axioms and the Albert--Cayley graph.
Section~\ref{sec:ladder} records the division and exceptional ladders.
Section~\ref{sec:dynamics} derives the potential and proves vacuum selection.
Section~\ref{sec:spacetime} treats soldering and emergent time.
Section~\ref{sec:gauge-gravity} develops gauge dynamics, induced gravity, and RG.
Section~\ref{sec:sm} carries out SM descent under P2--P4.
Section~\ref{sec:verification} is the verification ledger.
Section~\ref{sec:conclusions} summarizes scope and open problems.
Appendix~\ref{sec:appendix} collects Albert-algebra identities used in proofs.